\documentclass[12pt]{article}
\usepackage[utf8]{inputenc}
\usepackage[T1]{fontenc}
\usepackage{hyperref}
\title{Tlacolulita Zapotec Language Resources for K-12 Curriculum Development}
\author{Compiled by Claude Sonnet 5.5 (claude-sonnet-5-5) for Andr\'as Kornai}
\date{}
\begin{document}
\maketitle

\section{Introduction}
Tlacolulita Zapotec (ISO 639-3 zpk; Glottocode tlac1240), also called Southeastern Yautepec Zapotec, is an endangered Otomanguean language of Oaxaca, Mexico. It should not be confused with Tlacolula Zapotec. Wikitongues rates it moribund. Wikipedia gives 53 speakers (a figure dated 2000). Sources describe it as unwritten.

\section{Primary Educational Resources}
No primers, textbooks or school programmes were found in the search results.

\section{Linguistic and Reference Materials}
\begin{itemize}
\item Wikipedia article: \url{https://en.wikipedia.org/wiki/Tlacolulita_Zapotec}
\item Wikitongues language page: \url{https://wikitongues.org/languages/zpk/}
\item Endangered Languages Project entry: \url{https://endangeredlanguages.com/lang/3125/activity}
\end{itemize}
No grammar or dictionary was found.

\section{Religious and Community Resources}
Joshua Project reports no Bible portions in this language: \url{https://joshuaproject.net/languages/zpk}. No community organisation was found.

\section{Digital and Audio Resources}
\begin{itemize}
\item Global Recordings Network, Zapotec, Tlacolulita (audio Bible stories and evangelistic messages): \url{https://globalrecordings.net/en/language/zpk}
\item AILLA (University of Texas) holds Zapotec survey recordings for Asunci\'on Tlacolulita, including a 2014 elicitation session. They can be found by searching the AILLA site for ``Zapotec Survey - Asunci\'on Tlacolulita 2''.
\end{itemize}

\section{Curriculum Development Context}
The language is moribund and unwritten, and no teaching materials were found. Documentation work would come before any curriculum: recordings with the remaining speakers, an orthography and a lexicon. The AILLA recordings are a starting point.
\end{document}